O gerenciamento lógico da infraestrutura de rede compreende o planejamento, implementação e manutenção da configuração da camada 3 (roteamento) e camada 4(firewall).

\begin{itemize} 
    \item Implementação de NAT quando necessário; 
    \item Monitoramento de tráfego para identificação de gargalos ou comportamentos anômalos; 
    \item Revisão periódica das regras de firewall e políticas de acesso; 
    \item Definição de políticas de comunicação entre VLANs (no HP são as Zonas); 
    \item Criação e manutenção de regras de firewall para controle de acesso entre redes internas e externas; 
    \item Planejamento de endereçamento IP e organização de sub-redes; 
\end{itemize}


